\section{Notaci\'on}
\label{sec:notacion}

Este documento usa la convenci\'on de \cite{russell} para distinguir
distribuciones de probabilidades puntuales:

\begin{itemize}
  \item $\prob(a)$ es un \emph{n\'umero}: la probabilidad de que la variable
    tome el valor $a$.
  \item $\Prob(A)$ es un \emph{vector} de n\'umeros: la distribuci\'on de
    probabilidad sobre todos los valores del dominio de $A$. As\'i,
    $\Prob(A\given B)$ denota la tabla completa de probabilidades
    condicionales.
\end{itemize}

Para secuencias temporales se escribe
\begin{equation}
  \label{eq:notacion_secuencia}
  \begin{gathered}
    \Xr{m}{n} = \X{m},\,\X{m+1},\,\dots,\,\X{n},\\
    m,n\in\set{0,1,2,\dots},\quad m\le n .
  \end{gathered}
\end{equation}

S\'imbolos recurrentes:

\begin{center}
\small
\begin{tabular}{@{}ll@{}}
  \toprule
  S\'imbolo & Significado \\
  \midrule
  $\Tm$        & tiempo de muestreo \\
  $t_k=k\,\Tm$ & instante discreto $k$-\'esimo \\
  $\X{k}$      & estado (no observable) en $t_k$ \\
  $\Ev{k}$     & evidencia (observable) en $t_k$ \\
  $\ev{k}$     & valor concreto observado de $\Ev{k}$ \\
  $\al$        & constante de normalizaci\'on \\
  $\Verd,\Fal$ & valores verdadero y falso \\
  \bottomrule
\end{tabular}
\end{center}
